\documentclass[10pt,a4paper]{amsart}
\usepackage{amsmath,amssymb,amsthm,graphicx}
\usepackage{verbatim}
\newenvironment{reference}{\comment}{\endcomment}
%\newenvironment{reference}{}{}
\newenvironment{slogan}{\comment}{\endcomment}
\newenvironment{history}{\comment}{\endcomment}

% For commutative diagrams we use Xy-pic
\usepackage[all]{xy}

% We use 2cell for 2-commutative diagrams.
\xyoption{2cell}
% Theorem environments.
%
\theoremstyle{plain}
\newtheorem{theorem}[subsection]{Theorem}
\newtheorem{proposition}[subsection]{Proposition}
\newtheorem{lemma}[subsection]{Lemma}

\theoremstyle{definition}
\newtheorem{definition}[subsection]{Definition}
\newtheorem{example}[subsection]{Example}
\newtheorem{exercise}[subsection]{Exercise}
\newtheorem{situation}[subsection]{Situation}

\theoremstyle{remark}
\newtheorem{remark}[subsection]{Remark}
\newtheorem{remarks}[subsection]{Remarks}

\numberwithin{equation}{subsection}

% Macros
%
\def\lim{\mathop{\mathrm{lim}}\nolimits}
\def\colim{\mathop{\mathrm{colim}}\nolimits}
\def\Spec{\mathop{\mathrm{Spec}}}
\def\Hom{\mathop{\mathrm{Hom}}\nolimits}
\def\Ext{\mathop{\mathrm{Ext}}\nolimits}
\def\SheafHom{\mathop{\mathcal{H}\!\mathit{om}}\nolimits}
\def\SheafExt{\mathop{\mathcal{E}\!\mathit{xt}}\nolimits}
\def\Sch{\mathit{Sch}}
\def\Mor{\mathop{\mathrm{Mor}}\nolimits}
\def\Ob{\mathop{\mathrm{Ob}}\nolimits}
\def\Sh{\mathop{\mathit{Sh}}\nolimits}
\def\NL{\mathop{N\!L}\nolimits}
\def\CH{\mathop{\mathrm{CH}}\nolimits}
\def\proetale{{pro\text{-}\acute{e}tale}}
\def\etale{{\acute{e}tale}}
\def\QCoh{\mathit{QCoh}}
\def\Ker{\mathop{\mathrm{Ker}}}
\def\Im{\mathop{\mathrm{Im}}}
\def\Coker{\mathop{\mathrm{Coker}}}
\def\Coim{\mathop{\mathrm{Coim}}}

% Boxtimes
%
\DeclareMathSymbol{\boxtimes}{\mathbin}{AMSa}{"02}

%
% Macros for moduli stacks/spaces
%
\def\QCohstack{\mathcal{QC}\!\mathit{oh}}
\def\Cohstack{\mathcal{C}\!\mathit{oh}}
\def\Spacesstack{\mathcal{S}\!\mathit{paces}}
\def\Quotfunctor{\mathrm{Quot}}
\def\Hilbfunctor{\mathrm{Hilb}}
\def\Curvesstack{\mathcal{C}\!\mathit{urves}}
\def\Polarizedstack{\mathcal{P}\!\mathit{olarized}}
\def\Complexesstack{\mathcal{C}\!\mathit{omplexes}}
% \Pic is the operator that assigns to X its picard group, usage \Pic(X)
% \Picardstack_{X/B} denotes the Picard stack of X over B
% \Picardfunctor_{X/B} denotes the Picard functor of X over B
\def\Pic{\mathop{\mathrm{Pic}}\nolimits}
\def\Picardstack{\mathcal{P}\!\mathit{ic}}
\def\Picardfunctor{\mathrm{Pic}}
\def\Deformationcategory{\mathcal{D}\!\mathit{ef}}
\usepackage{cleveref}
\crefname{originallemma}{original lemma}{original lemmas}
\Crefname{originallemma}{Original lemma}{Original lemmas}
\makeatletter
\newlabel{lemma-gabriel-popescu}{{0F5V}{source}{}{sourcepointer-1}{}}
\newlabel{lemma-gabriel-popescu@cref}{{[originallemma][0][]0F5V}{[1][0][]source}{}{}{}}
\makeatother
\crefname{originallemma}{original lemma}{original lemmas}
\Crefname{originallemma}{Original lemma}{Original lemmas}
\makeatletter
\newlabel{derived-lemma-brown}{{0A8F}{source}{}{sourcepointer-2}{}}
\newlabel{derived-lemma-brown@cref}{{[originallemma][0][]0A8F}{[1][0][]source}{}{}{}}
\makeatother
\crefname{originallemma}{original lemma}{original lemmas}
\Crefname{originallemma}{Original lemma}{Original lemmas}
\makeatletter
\newlabel{derived-lemma-split}{{05QT}{source}{}{sourcepointer-3}{}}
\newlabel{derived-lemma-split@cref}{{[originallemma][0][]05QT}{[1][0][]source}{}{}{}}
\makeatother
\crefname{originallemma}{original lemma}{original lemmas}
\Crefname{originallemma}{Original lemma}{Original lemmas}
\makeatletter
\newlabel{derived-lemma-projectors-have-images-triangulated}{{05QW}{source}{}{sourcepointer-4}{}}
\newlabel{derived-lemma-projectors-have-images-triangulated@cref}{{[originallemma][0][]05QW}{[1][0][]source}{}{}{}}
\makeatother
\title{Brown representability for functors on the derived category}
\author{The Stacks Project Authors}
\date{stacks-project:injectives.tex:local-snapshot-20261012; GNU FDL 1.2 or later}
\begin{document}
\maketitle
\noindent\textit{Editorial note.} All delivered mathematical text is the author's own, copied from the pinned Stacks Project chapter file sites-cohomology.tex; no statement or proof step was rewritten, reordered or removed.\par\medskip
\noindent\textit{Editorial note.} The proof is the author's complete argument; the comparison is displayed as a chain of identifications ending in a Hom set, the splitting is functorial in the test object, and the final step names the Karoubian property used.\par\medskip
\noindent\textit{Editorial note.} Every labelled reference occurring inside the counted spans is carried in assembly.external\_references as a non-load-bearing navigation pointer (same-chapter ones at their definition site, cross-chapter ones at the invocation inside the counted span), so no load-bearing mathematics is replaced by a stub. No bibliography citation occurs inside the counted spans.\par\medskip
\noindent\textit{Editorial note.} The same-name check for this card scanned every staged Stacks chapter of this repository -- all sources/stacks-*/pinned-original-source.tex files, the card's own chapter excluded -- and the labels found are recorded in dedup.same\_name\_cross\_chapter on the card.\par\medskip
\noindent\textit{Editorial note.} The preamble is the chapter's own retained minimal preamble (sources/stacks-sites-cohomology/retained-preamble.tex), a verbatim excerpt of the native preamble: the Xy-pic setup and the native theorem declarations with the mathematical macros.\par\medskip
\noindent\textit{Allowed standard prerequisites (editorial).}\begin{itemize}
\item Brown representability for the abelian category. Same chapter, Lemma lemma-grothendieck-brown. Supplies representability for functors on the abelian category commuting with colimits, which the author upgrades to the derived category setting.
\item The derived category is Karoubian. Derived Categories, Lemma derived-lemma-projectors-have-images-triangulated. Supplies that idempotents split in the derived category, which turns the representation up to a direct summand into an honest representation.
\item Cohomological functors and their representability. Same chapter, the discussion of cohomological functors on the derived category and the functors they induce. Supplies the passage from a functor on the derived category to the collection of functors on the abelian category used in the argument.
\end{itemize}\medskip
\ifdefined\hypertarget\hypertarget{sourcepointer-1}{}\fi
\noindent\textit{Source reference.} Original lemma 0F5V: sources/stacks-injectives/pinned-original-source.tex, source lines 2922--2922. Navigation pointer invoked at proof line 2983; the counted span contains the pointer only and the cited result is not reproduced inside it, so the pointer is registered as a non-load-bearing navigation stub in the established form for Stacks material. No bibliography citation is involved.\par\medskip
\ifdefined\hypertarget\hypertarget{sourcepointer-2}{}\fi
\noindent\textit{Source reference.} Original lemma 0A8F: sources/stacks-derived/pinned-original-source.tex, source lines 11546--11546. Navigation pointer invoked at proof line 2989; the counted span contains the pointer only and the cited result is not reproduced inside it, so the pointer is registered as a non-load-bearing navigation stub in the established form for Stacks material. No bibliography citation is involved.\par\medskip
\ifdefined\hypertarget\hypertarget{sourcepointer-3}{}\fi
\noindent\textit{Source reference.} Original lemma 05QT: sources/stacks-derived/pinned-original-source.tex, source lines 917--917. Navigation pointer invoked at proof line 2998; the counted span contains the pointer only and the cited result is not reproduced inside it, so the pointer is registered as a non-load-bearing navigation stub in the established form for Stacks material. No bibliography citation is involved.\par\medskip
\ifdefined\hypertarget\hypertarget{sourcepointer-4}{}\fi
\noindent\textit{Source reference.} Original lemma 05QW: sources/stacks-derived/pinned-original-source.tex, source lines 1023--1023. Navigation pointer invoked at proof line 3017; the counted span contains the pointer only and the cited result is not reproduced inside it, so the pointer is registered as a non-load-bearing navigation stub in the established form for Stacks material. No bibliography citation is involved.\par\medskip
\begin{lemma}
\label{lemma-brown}
Let $\mathcal{A}$ be a Grothendieck abelian category.
Let $H : D(\mathcal{A}) \to \textit{Ab}$ be a contravariant
cohomological functor which transforms direct sums into products.
Then $H$ is representable.
\end{lemma}

\begin{proof}
Let $R, F, G, RG$ be as in Lemma \ref{lemma-gabriel-popescu}
and consider the functor $H \circ F : D(\text{Mod}_R) \to \textit{Ab}$.
Observe that since $F$ is a left adjoint it sends direct sums to
direct sums and hence $H \circ F$ transforms direct sums into products.
On the other hand, the derived category $D(\text{Mod}_R)$ is
generated by a single compact object, namely $R$.
By Derived Categories, Lemma \ref{derived-lemma-brown}
we see that $H \circ F$ is representable, say by $L \in D(\text{Mod}_R)$.
Choose a distinguished triangle
$$
M \to L \to RG(F(L)) \to M[1]
$$
in $D(\text{Mod}_R)$. Then $F(M) = 0$ because $F \circ RG = \text{id}$.
Hence $H(F(M)) = 0$ hence $\Hom(M, L) = 0$.
It follows that $L \to RG(F(L))$ is the inclusion of a direct summand, see
Derived Categories, Lemma \ref{derived-lemma-split}.
For $A$ in $D(\mathcal{A})$ we obtain
\begin{align*}
H(A)
& =
H(F(RG(A)) \\
& =
\Hom(RG(A), L) \\
& \to
\Hom(RG(A), RG(F(L))) \\
& =
\Hom(F(RG(A)), F(L)) \\
& =
\Hom(A, F(L))
\end{align*}
where the arrow has a left inverse functorial in $A$. In other words, we find
that $H$ is the direct summand of a representable functor.
Since $D(\mathcal{A})$ is Karoubian
(Derived Categories, Lemma
\ref{derived-lemma-projectors-have-images-triangulated}) we conclude.
\end{proof}

\end{document}
